% =============================================================================
\chapter{O aplicativo cliente}
\label{cap:cliente}
% =============================================================================

O aplicativo é o que o aluno vê. Escrito em Python, com Tkinter para as janelas
e Matplotlib para os gráficos, roda em qualquer PC do laboratório e não precisa
de nenhum software da Intel. As dependências são o NumPy e o Matplotlib; o SciPy
é opcional e só é usado pelo projeto de filtros IIR elípticos.

\section{Abrir o aplicativo}

Na instalação de turma (capítulo~\ref{cap:laboratorio}):

\begin{center}
\ttfamily cd /opt/morphe/Python \&\& python3 morphe\_app.py
\end{center}

A janela principal abre já escolhendo a placa: entre as do \arq{placas.conf} e
as que a instalação conhece, fica com a livre e menos usada, e varre a rede em
segundo plano atrás de placas fora da lista (capítulo~\ref{cap:placas}).

O painel \textbf{Placas do laboratório}, no topo, mostra uma linha por placa,
atualizada a cada 4\,s: um ponto de cor (verde livre, âmbar ocupada, cinza sem
resposta, vermelho com a FPGA não preparada), o nome, o IP e o que a placa está
fazendo --- ``ocupada: convolução há 0,1\,s · 3 computadores no último
minuto''. Em cima, um resumo (``2 placas · 1 livre · 1 ocupada'') e o botão
\emph{Procurar de novo}. A placa em uso é marcada com ``você está aqui''; as
outras têm um botão \emph{Usar}, para trocar de placa com um clique. Sem placa
nenhuma, o painel diz isso numa linha, sem janela de erro: o gerador de sinais,
o comparador e os projetistas de filtro funcionam sem placa. O endereço e a
porta, para conectar à mão, ficam recolhidos em \emph{Conexão manual}.

\section{As janelas}

A janela principal é um menu de botões, um por ferramenta:

\begin{description}
  \item[Gerador de Sinais] monta um sinal e o mostra, sem placa. Serve para
    explorar e para salvar sinais em arquivo.
  \item[Aquisição (ADC da placa)] lê tensões e captura sinais pela entrada
    analógica da placa (seção~\ref{sec:janela-aquisicao}).
  \item[Convolução] dois sinais, \(x[n]\) e \(h[n]\), convoluídos na FPGA. Mostra
    as entradas, a saída da placa, a referência do NumPy e o erro entre elas.
  \item[FFT] um sinal, transformado na FPGA; mostra magnitude e fase, compara com
    o NumPy e permite voltar ao tempo pela IFFT da própria placa.
  \item[IFFT] um espectro lido de arquivo (\arq{.mrph}, \arq{.npy} complexo ou
    texto com colunas real e imaginária), transformado de volta na FPGA. Serve
    para testar a inversa isolada do erro da direta.
  \item[Filtro FIR] um sinal montado por superposição de componentes, filtrado
    na FPGA por um FIR que vem do projetista ou de um arquivo de coeficientes.
  \item[Filtro IIR] a mesma ideia, com um IIR em seções de segunda ordem. O
    gráfico do meio é a resposta em frequência do projeto e do filtro
    quantizado. A saída da placa é conferida na hora, bit a bit, contra o
    modelo em ponto fixo, e a barra de status diz se bateu.
  \item[Comparar .mrph] abre um pacote salvo (seção~\ref{sec:mrph}) e refaz a
    conta no computador, sem placa.
\end{description}

Todas as janelas que usam a placa têm a mesma estrutura: um painel para o sinal
de entrada, um botão que envia à FPGA, os gráficos (com barra de ferramentas
para ampliar, deslocar e destacar o gráfico numa janela própria) e uma barra de
status que diz o que aconteceu, inclusive o tempo que a placa levou.

Existe também uma janela de espectrograma, pronta e validada, que ficou fora do
menu por decisão do supervisor; as linhas que a ligam estão comentadas no
\arq{morphe_app.py}, com a marca \texttt{ESPECTROGRAMA}.

\subsection*{De onde vem o sinal}

O painel de sinal de todas as janelas oferece sete tipos:

\begin{itemize}
  \item cinco gerados: degrau, impulso, senoide, exponencial e retangular, com
    os seus parâmetros e o número de amostras \(N\);
  \item \textbf{Arquivo}: um sinal qualquer, lido de \arq{.csv}, \arq{.txt},
    \arq{.npy} ou \arq{.wav}. Texto com uma coluna (as amostras) ou duas
    (índice ou tempo, amostra), com vírgula, ponto e vírgula, tabulação ou
    espaços; aceita a vírgula decimal de planilhas em português. Do
    \arq{.wav} vem também a taxa de amostragem. O leitor está em
    \arq{Python/sinal_arquivo.py};
  \item \textbf{Captura do ADC}: a última captura feita na janela de aquisição.
\end{itemize}

Com os dois últimos o tamanho do sinal é o do arquivo ou da captura, e pode
passar muito de 1024 amostras: as janelas então processam por blocos
(seção~\ref{sec:blocos}). Quando o sinal vem com uma taxa de amostragem e o
filtro foi projetado com outra, as janelas do FIR e do IIR avisam.
\arq{exemplos/sinais/} tem arquivos de exemplo para cada janela, com um
\arq{README.md}.

\subsection*{Os projetistas de filtros}

\paragraph{FIR.} O projetista (\arq{fir_design.py}) tem dois modos. \emph{Por
especificação}, o padrão: o aluno dá o tipo (passa-baixa, passa-alta,
passa-banda, rejeita-banda), a borda da banda passante, a largura da transição,
o \emph{ripple} e a atenuação; o programa escolhe a janela (retangular, Hann,
Hamming, Blackman ou Kaiser) e calcula quantos coeficientes são necessários. É
o método do \arq{dsp_fir_filter.m} do Prof.\ Sanca, portado para Python. No modo
\emph{avançado}, o aluno escolhe o número de coeficientes e a janela. Nos dois,
o projetista mostra o que o filtro entrega \emph{depois} da quantização em
Q15.16. Validado na placa em 22/09/2026: 63,0\,dB de rejeição medida contra
63,3\,dB da referência.

\paragraph{IIR.} O projetista (\arq{iir_design.py}) segue o
\arq{dsp_iir_filter.m} do Prof.\ Sanca: Butterworth, Chebyshev~I, Chebyshev~II
e elíptico, nos quatro tipos. A diferença deliberada é a saída: em vez de um
polinômio de ordem \(N\), ela é uma cascata de seções de segunda ordem, pelos
motivos da seção~\ref{sec:acel-iir}. O Butterworth e os dois Chebyshev são
calculados só com NumPy; o elíptico usa o SciPy e só aparece se ele estiver
instalado. Quatro erros encontrados no \arq{.m} original estão corrigidos no
módulo Python e documentados nele.

\section{A janela de aquisição}
\label{sec:janela-aquisicao}

A janela ``Aquisição (ADC da placa)'' depende de a placa ter o bitstream com o
ADC; se não tiver, a própria mensagem de recusa do servidor aparece na barra de
status (seção~\ref{sec:marca}). Ela faz duas coisas.

\paragraph{Voltímetro.} Lê a tensão média de 0,1\,s na entrada escolhida, com
desvio, mínimo e máximo, e pode repetir sozinho para acompanhar um ajuste. Não
há hardware separado para isso: é uma captura de 2000 amostras a 20\,kHz. O
intervalo de 0,1\,s contém um número inteiro de ciclos de 50\,Hz e de 60\,Hz,
então a média cancela o zumbido da rede elétrica.

\paragraph{Captura.} O aluno escolhe a entrada (um dos oito canais contra o
terra, de 0 a 4,095\,V, ou um dos quatro pares diferenciais, de
\(\pm 2{,}048\)\,V), a taxa de amostragem (de 1 a 200\,kHz) e o número de
amostras:

\begin{itemize}
  \item até 32\,768, a captura é de uma vez, na memória da FPGA;
  \item acima disso, ou com 0 amostras (sem limite, até clicar em
    \emph{Parar}), é contínua: a placa envia em blocos enquanto captura
    (seção~\ref{sec:adc-continuo}).
\end{itemize}

O gráfico mostra o sinal em volts; acima de 200\,mil pontos, a envoltória. Para
uma senoide, a janela mede frequência, SINAD e ENOB, que é o jeito de conferir o
conversor com um gerador de funções (seção~\ref{sec:adc}). A captura pode ser
salva em \arq{.csv}, \arq{.npy} ou \arq{.wav} e passa a estar disponível como o
tipo ``Captura do ADC'' em todas as outras janelas. As contas --- a palavra de
configuração, a taxa real, a conversão para volts, as métricas e a gravação ---
estão em \arq{Python/aquisicao.py}, separadas da janela.

A taxa real é \(50\,\text{MHz}/\text{divisor}\), com divisor inteiro. Ela é
exata para as taxas que dividem 50\,MHz; as outras são arredondadas, e a janela
mostra a taxa real usada (44,1\,kHz, por exemplo, vira 44\,091,7\,Hz).

\section{Os módulos}

As janelas são finas; o que importa está nos módulos que elas usam.

\begin{table}[htbp]
  \centering
  \small
  \caption{Os módulos do cliente, em \arq{Python/}.}
  \label{tab:modulos}
  \begin{tabularx}{\textwidth}{@{}lX@{}}
    \toprule
    \textbf{Módulo} & \textbf{O que faz} \\
    \midrule
    \arq{morphe_app.py} & a janela principal e o menu \\
    \arq{*_window.py} & uma janela por ferramenta \\
    \arq{signal_panel.py} & o painel de sinal comum a todas as janelas \\
    \arq{tcp_panel.py} & o painel ``Placas do laboratório'': descoberta, escolha e estado das placas \\
    \arq{morphe_protocol.py} & o protocolo (capítulo~\ref{cap:protocolo}):
      monta e decodifica as mensagens, abre as conexões, varre a rede \\
    \arq{morphe_config.py} & os limites do hardware, espelho do
      \arq{C/morphe_config.h} \\
    \arq{dsp_core.py} & geração de sinais, conversão para ponto fixo, a
      referência em ponto flutuante e a leitura e escrita de \arq{.mrph} \\
    \arq{blocos.py} & o processamento de sinais maiores que o hardware
      (seção~\ref{sec:blocos}) \\
    \arq{sinal_arquivo.py} & a leitura de sinais de arquivo \\
    \arq{aquisicao.py} & as contas do ADC \\
    \arq{fir_design.py}, \arq{iir_design.py} & os projetistas, e o modelo em
      ponto fixo do IIR que reproduz a placa bit a bit \\
    \arq{morphe_theme.py}, \arq{plot_toolbar.py}, \arq{popout_helper.py}
      & aparência, barra dos gráficos, gráficos em janela própria \\
    \arq{morphe_export.py} & exporta uma janela aberta para SVG vetorial, para
      figuras de artigos e deste manual \\
    \bottomrule
  \end{tabularx}
\end{table}

\section{O pacote \texttt{.mrph} e o comparador}
\label{sec:mrph}

O \arq{.mrph} é o formato de arquivo do Morphe: texto puro, legível em
qualquer editor. Um arquivo simples tem um cabeçalho com linhas iniciadas por
\texttt{\#} (data, tipo numérico, tamanho, taxa de amostragem, descrição) e uma
tabela de valores, com duas colunas para sinais reais (\(n\), valor) e três para
complexos (\(k\), real, imaginário). Um \emph{pacote} reúne vários sinais num
arquivo só, em seções iniciadas por \texttt{\#\# section:}, cada uma com o seu
cabeçalho: numa convolução, \(x\), \(h\) e \(y\); numa FFT, \(x\) e \(X\).

Os pacotes vêm de dois lugares: as janelas salvam o que o aluno fez, e o
servidor grava um a cada operação na placa (seção~\ref{sec:marca} e seguinte).
O comparador abre qualquer um deles, identifica a operação, refaz a conta em
NumPy (no IIR, no modelo em ponto fixo) sobre a mesma entrada e mostra a
diferença. É o caminho para estudar um resultado depois, sem placa, ou para
investigar um resultado estranho que um aluno trouxe. \arq{exemplos/} tem
pacotes reais produzidos pela placa.

\section{Sinais maiores que o hardware: processamento em blocos}
\label{sec:blocos}

Os aceleradores processam no máximo 1024 amostras por vez. Um arquivo de áudio
de um segundo tem 44\,100; uma captura do ADC pode ter milhões. O
\arq{Python/blocos.py} divide o sinal em blocos, manda cada um à placa pelo
mesmo caminho que as janelas usam para um bloco só, e junta os resultados. São
três técnicas, uma por tipo de operação, e cada uma tem um argumento de por que
o resultado é o mesmo que uma operação única daria.

Todas as funções de bloco recebem a operação de \emph{um} bloco como parâmetro:
na placa, as classes \texttt{ConvPlaca}, \texttt{FirPlaca}, \texttt{FftPlaca},
\texttt{IfftPlaca} e \texttt{IirPlaca}; nos testes sem placa, o
\texttt{np.convolve}, o \texttt{np.fft.fft} e o modelo em ponto fixo. Isso
permite conferir a montagem dos blocos separada do hardware.

\subsection{Convolução e FIR: \emph{overlap-add}}

A convolução é linear e invariante no tempo. Cortando \(x\) em blocos \(x_i\) de
\(L\) amostras,
\[
  x[n] = \sum_i x_i[n - iL]
  \quad\Longrightarrow\quad
  (x * h)[n] = \sum_i (x_i * h)[n - iL].
\]
Cada \(x_i * h\) é uma convolução que cabe na placa, com \(L + M - 1\) amostras
de saída; as caudas de \(M - 1\) amostras de blocos vizinhos se sobrepõem e são
somadas (Figura~\ref{fig:overlap-add}). Se \(h\) também passa de 1024
amostras, é cortado do mesmo jeito, e cada par de pedaços vira uma requisição.
O resultado é exatamente a convolução linear: a única diferença para uma
operação única é a quantização em Q15.16, que já existia.

\begin{figure}[htbp]
  \centering
  \figura{overlap-add}
  \caption{\emph{Overlap-add}: cada bloco de \(L\) amostras é convoluído na
    placa; as caudas de \(M-1\) amostras se somam.}
  \label{fig:overlap-add}
\end{figure}

\subsection{FFT e IFFT: o algoritmo de quatro passos}

Uma DFT de \(N = M \cdot 1024\) pontos se decompõe em DFTs de 1024 e DFTs de
\(M\) pontos, ligadas por fatores de giro. É a decomposição de Cooley--Tukey
com \(N_1 = M\) e \(N_2 = 1024\). Com \(n = n_1 + M n_2\) e
\(k = 1024\,k_1 + k_2\):
\[
  X[1024\,k_1 + k_2] = \sum_{n_1=0}^{M-1} W_M^{n_1 k_1}\;
    W_N^{n_1 k_2}\;
    \underbrace{\mathrm{FFT}_{1024}\{x[n_1 + M n_2]\}[k_2]}_{\text{na placa}}.
\]
Em quatro passos (Figura~\ref{fig:quatro-passos}):

\begin{enumerate}
  \item a placa faz \(M\) FFTs de 1024 pontos, uma por subsequência
    \(x[n_1], x[n_1 + M], x[n_1 + 2M], \ldots\);
  \item cada resultado é multiplicado pelos fatores de giro \(W_N^{n_1 k_2}\);
  \item o PC faz 1024 DFTs de \(M\) pontos, uma por coluna \(k_2\);
  \item a matriz resultante, lida por linhas, é \(X[k]\) na ordem natural.
\end{enumerate}

\begin{figure}[htbp]
  \centering
  \figura{quatro-passos}
  \caption{FFT de \(N = M\cdot 1024\) pontos pelo algoritmo de quatro passos:
    \(M\) FFTs de 1024 na placa, fatores de giro e DFTs de \(M\) pontos no PC.}
  \label{fig:quatro-passos}
\end{figure}

Não é uma aproximação: é a mesma DFT, com a conta dividida. A precisão é a das
FFTs da placa, cada uma normalizada por conta própria, mais a do ponto
flutuante de 64 bits do PC, desprezível perto dela. O sinal é completado com
zeros até o múltiplo de 1024 acima. A IFFT longa é a mesma decomposição no
sentido inverso, e por isso exige que o espectro já tenha um número de pontos
múltiplo de 1024.

A escolha foi deliberada: um espectrograma (FFTs de quadros sobrepostos) também
processaria sinais longos, mas mudaria o significado da janela da FFT, que é
mostrar \emph{a} transformada do sinal.

\subsection{IIR: blocos com aquecimento}

O IIR não se decompõe assim. A saída em \(n\) depende de toda a entrada
anterior, e o bloco da placa sempre começa com o estado zerado. A solução é a
usual: cada bloco começa \(A\) amostras \emph{antes} do trecho que interessa, e
as primeiras \(A\) saídas são descartadas (Figura~\ref{fig:iir-blocos}). Se o
aquecimento for longo o bastante para o transitório do estado zero morrer, o
trecho guardado coincide com o do filtro rodando sem parar.

\begin{figure}[htbp]
  \centering
  \figura{iir-blocos}
  \caption{IIR por blocos: cada bloco começa algumas amostras antes
    (\emph{aquecimento}) e essa parte da saída é descartada.}
  \label{fig:iir-blocos}
\end{figure}

O aquecimento é calculado pelo polo de maior módulo \(|p|\): o transitório
decai como \(|p|^n\), e o critério é ele cair abaixo de meio passo do Q15.16,
\[
  A = \left\lceil 1{,}5 \cdot \frac{\ln 2^{-17}}{\ln |p|} \right\rceil
      + 8 \cdot (\text{número de seções}),
\]
com a folga de 50\,\% e o termo por seção cobrindo os polos repetidos. Um filtro
com polos muito perto do círculo unitário pode precisar de um aquecimento maior
que o próprio bloco; nesse caso o cliente recusa com uma mensagem que explica
por quê.

Para a conferência bit a bit da janela continuar valendo, o modelo em Python
roda com \emph{os mesmos blocos}: placa e modelo batem exatamente. A distância
entre os blocos e o filtro contínuo é medida à parte e mostrada em LSB. Em ponto
fixo ela nem sempre chega a zero, porque o arredondamento na realimentação pode
fazer estados iniciais diferentes convergirem para bits diferentes: medido em
23/09/2026, 0\,LSB num Butterworth de ordem 4 com corte em 1\,kHz e 4 a 8\,LSB
num de ordem 2 com corte em 100\,Hz, ambos a 8\,kHz.

\subsection{O custo}

Cada bloco é uma requisição, e a placa atende cerca de 8 por segundo
(capítulo~\ref{cap:placas}). Uma convolução de 5000 amostras por um filtro de 67
coeficientes são 5 requisições; um FIR sobre 5\,s de áudio a 44,1\,kHz são 216,
perto de meio minuto. Enquanto isso, a barra de status das janelas mostra o
bloco em curso (``Bloco 3 de 5 na FPGA'').
